\section{Conclusion}

This article is motivated by the knot selection problem in the multivariate spline regression. %We take into account the complexity of the model space and specify a prior on the number of knots to hedge the complexity. 
The proposed method estimates the number and location of spline knots simultaneously and accurately. %regardless of the true number of knots and the sample size. 
It provides a mechanism for easy interpretation and function fitting with jumping discontinuity. The trade-off for advantages is the additional computational cost, attributed to RJMCMC. Nevertheless, it will not cause a serious problem in small or moderately large datasets.

There are several potential jobs in the future. We intend to establish a theoretical framework of the algorithm including the consistency and convergence rate. The primary challenge is how to define an appropriate distance between parameters with varying dimensions. Furthermore, the domain space is required to be a cube. We aim for an extension in the general domain. A possible solution is substituting a flexible multivariate spline for the tensor product spline.